\section*{Chapitre \ref{ch:b3:representations} --- Représentations
des groupes finis}

\begin{solution}{exo:b3:representations:1}
(a) $\Z/n\Z$ est abélien~: toutes les \omterm{def:b3:representations:rep}{irréductibles} sont de \omterm{def:b3:galois:extension}{degré}
$1$ (\cref{prop:b3:representations:onedim}), c'est-à-dire des
morphismes $\chi \colon \Z/n\Z \to \C^\times$, déterminés par
$\chi(\bar 1) = \omega$ avec $\omega^n = 1$~: les $n$ \omterm{def:b3:representations:character}{caractères}
$\chi_k(\bar m) = \eu^{2\iu\pi km/n}$. Pour $n = 4$ (classes $=$
éléments $\bar0,\bar1,\bar2,\bar3$)~:
\[
\begin{array}{c|cccc}
& \bar0 & \bar1 & \bar2 & \bar3\\
\hline
\chi_0 & 1 & 1 & 1 & 1\\
\chi_1 & 1 & \iu & -1 & -\iu\\
\chi_2 & 1 & -1 & 1 & -1\\
\chi_3 & 1 & -\iu & -1 & \iu
\end{array}
\]
(b) Lignes~: $\langle\chi_k,\chi_l\rangle = \frac14\sum_m
\eu^{2\iu\pi(l-k)m/4} = \delta_{kl}$ (somme géométrique)~; colonnes
de même. La matrice $(\eu^{2\iu\pi km/n})_{k,m}$ est la matrice de
la \emph{transformée de Fourier discrète} --- le même filtre en
racines de l'unité que dans le chapitre des fonctions génératrices
du volume de L2~; l'\omterm{thm:b3:representations:orthogonality}{orthogonalité des caractères} généralise la
formule d'inversion de la TFD.
\end{solution}

\begin{solution}{exo:b3:representations:2}
(a) La matrice de $\rho(g)$ dans la base $(e_x)$ est une matrice
de permutation, de trace le nombre de $x$ avec $g\cdot x = x$.
Alors
\[
\langle\mathbf 1, \chi\rangle = \frac1{\abs G}\sum_g
\abs{\operatorname{Fix}(g)} = \#\{\text{orbites}\}
\]
par le lemme de dénombrement de Burnside (\cref{exo:b3:groups:5})
--- de façon équivalente, on calcule la multiplicité de la
\omterm{def:b3:representations:rep}{représentation} triviale, dont l'espace isotypique est l'espace
des vecteurs $G$-invariants, de dimension le nombre d'\omterm{def:b3:groups:action}{orbites}
(un indicateur par \omterm{def:b3:groups:action}{orbite}).

(b) Comme $\operatorname{Fix}_{X\times X}(g) =
\operatorname{Fix}_X(g)^2$, la partie (a) appliquée à $X \times
X$ donne $\langle\chi,\chi\rangle = \frac1{\abs
G}\sum\abs{\operatorname{Fix}(g)}^2 = \#\{\text{orbites sur }
X\times X\}$ ($\chi$ est réel). En écrivant $\chi = \mathbf 1 +
\psi$~: $\langle\mathbf1,\chi\rangle = 1$ (transitivité), donc
$\langle\psi,\psi\rangle = \langle\chi,\chi\rangle - 1$. L'\omterm{def:b3:groups:action}{action}
sur $X\times X$ a la diagonale comme une \omterm{def:b3:groups:action}{orbite}~; il y a
exactement une autre \omterm{def:b3:groups:action}{orbite} ssi $G$ est transitif sur les paires
distinctes~: $\langle\psi,\psi\rangle = 1$ ssi $2$-transitive
(\cref{cor:b3:representations:multiplicity}).

(c) $S_n$ est $2$-transitive sur $\{1,\dots,n\}$ (envoyer toute
paire distincte n'importe où)~: $\psi = \chi_{\mathrm{std}}$ est
\omterm{def:b3:representations:rep}{irréductible}.
\end{solution}

\begin{solution}{exo:b3:representations:3}
$S_3$~: trois classes, $\sum n_i^2 = 6 = 1 + 1 + 4$~; les deux
\omterm{def:b3:representations:character}{caractères} de \omterm{def:b3:galois:extension}{degré} $1$ sont $\mathbf 1, \varepsilon$ ($[S_3 :
A_3] = 2$)~; la troisième ligne $(2, a, b)$ suit de
l'orthogonalité des colonnes avec la colonne de $e$~: $1 - 1 +
2a = 0$ et $1 + 1 + 2b = 0$~: $a = 0$, $b = -1$ --- la table de
\cref{ex:b3:representations:s3}.

Vérifications $S_4$ (lignes)~: $\langle\chi_{\mathrm{std}},
\varepsilon\chi_{\mathrm{std}}\rangle = \frac1{24}(9 - 6 + 3 + 0
- 6) = 0$~; $\langle\chi_2,\chi_2\rangle = \frac1{24}(4 + 0 + 12
+ 8 + 0) = 1$. Colonnes~: $e$ contre $(1\,2)$~: $1 - 1 + 0 + 3 -
3 = 0$~; $(1\,2)$ contre elle-même~: $1 + 1 + 0 + 1 + 1 = 4 =
\abs{Z_{S_4}((1\,2))} = 24/6$.

\omterm{def:b3:representations:character}{Caractère} de permutation sur $4$ points~: $(4, 2, 0, 1, 0) =
\mathbf 1 + \chi_{\mathrm{std}}$ (comptages de points fixes~;
soustraire la ligne du haut). Pour $\chi_{\mathrm{std}}^2 =
(9, 1, 1, 0, 1)$~:
\[
\langle\mathbf1,\cdot\rangle = \tfrac{9 + 6 + 3 + 0 + 6}{24} =
1,\quad
\langle\varepsilon,\cdot\rangle = 0,\quad
\langle\chi_2,\cdot\rangle = \tfrac{18 + 6}{24} = 1,\quad
\langle\chi_{\mathrm{std}},\cdot\rangle = 1,\quad
\langle\varepsilon\chi_{\mathrm{std}},\cdot\rangle = 1:
\]
$\chi_{\mathrm{std}}^2 = \mathbf 1 + \chi_2 +
\chi_{\mathrm{std}} + \varepsilon\chi_{\mathrm{std}}$
(dimensions~: $9 = 1 + 2 + 3 + 3$).
\end{solution}

\begin{solution}{exo:b3:representations:4}
Les deux groupes ont cinq classes et le motif de \omterm{def:b3:galois:extension}{degrés}
$(1,1,1,1,2)$ (quatre \omterm{def:b3:galois:extension}{degrés} $1$ issus de l'abélianisé $\cong
(\Z/2\Z)^2$, puis $\sum n_i^2 = 8$). En ordonnant les classes
$e$, $z$ (l'involution centrale~: $r^2$, resp.\ $-1$), et les
trois classes à deux éléments~:
\[
\begin{array}{c|ccccc}
& e & z & C_1 & C_2 & C_3\\
\hline
\chi^{(1)} & 1 & 1 & 1 & 1 & 1\\
\chi^{(2)} & 1 & 1 & 1 & -1 & -1\\
\chi^{(3)} & 1 & 1 & -1 & 1 & -1\\
\chi^{(4)} & 1 & 1 & -1 & -1 & 1\\
\chi^{(5)} & 2 & -2 & 0 & 0 & 0
\end{array}
\]
(la dernière ligne par orthogonalité des colonnes). Tables
identiques pour $D_4$ et $Q_8$, qui ne sont pas isomorphes
(\cref{pb:b3:groups:1}). La table \emph{capture} bien~: $\abs
G$, tailles des classes, le \omterm{ex:b3:groups:actions}{centre} ($\{g : \abs{\chi_i(g)} =
n_i\ \forall i\}$~: ordre $2$ dans les deux cas), l'abélianisé,
tout le treillis des sous-groupes normaux (noyaux et
intersections, \cref{exo:b3:representations:6}). Elle échoue à
capturer les ordres d'éléments~: $D_4$ a cinq involutions, $Q_8$
en a une --- le type d'isomorphisme est vraiment plus fin que
la \omterm{def:b3:representations:table}{table de caractères}.
\end{solution}

\begin{solution}{exo:b3:representations:5}
(a) Dans $S_4$, le centralisateur de $(1\,2\,3)$ a ordre $24/8
= 3$~: c'est $\langle(1\,2\,3)\rangle \subseteq A_4$. Donc
$Z_{A_4}((1\,2\,3))$ a ordre $3$ et la classe de $A_4$ a $12/3
= 4$ éléments~: les huit $3$-cycles se scindent en deux classes
de $A_4$ (représentées par $(1\,2\,3)$ et son inverse).

(b) $A_4/V \cong \Z/3\Z$ donne trois \omterm{def:b3:representations:character}{caractères} de \omterm{def:b3:galois:extension}{degré} $1$
($\omega = \eu^{2\iu\pi/3}$~; les classes de $3$-cycles s'envoient
sur $\bar1, \bar2$)~; la restriction de $\chi_{\mathrm{std}}$
reste \omterm{def:b3:representations:rep}{irréductible} ($\langle\chi,\chi\rangle = \frac1{12}(9 + 3
\cdot 1 + 0 + 0) = 1$)~:
\[
\begin{array}{c|cccc}
A_4 & e\,[1] & (1\,2)(3\,4)\,[3] & (1\,2\,3)\,[4] &
(1\,3\,2)\,[4]\\
\hline
\mathbf 1 & 1 & 1 & 1 & 1\\
\chi_\omega & 1 & 1 & \omega & \omega^2\\
\chi_{\bar\omega} & 1 & 1 & \omega^2 & \omega\\
\chi_3 & 3 & -1 & 0 & 0
\end{array}
\]
(c) Noyaux~: $\ker\chi_\omega = \ker\chi_{\bar\omega} = V$~;
$\ker\chi_3 = \{e\}$ (aucune autre entrée n'a \omterm{def:b3:modules:module}{module} $3$). Les
sous-groupes normaux sont les intersections de noyaux
(\cref{exo:b3:representations:6})~: $\{e\}$, $V$, $A_4$ --- en
particulier $A_4$ n'a pas de \omterm{def:b3:groups:normal}{sous-groupe normal} d'ordre $2$ ni
d'indice $2$.
\end{solution}

\begin{solution}{exo:b3:representations:6}
(a) Si $\chi(g) = \chi(e) = n$~: l'égalité dans $\abs{\chi(g)}
\leq n$ force $\rho(g) = \lambda\,\mathrm{id}$
(\cref{prop:b3:representations:charbasics}) avec $n\lambda =
n$~: $\rho(g) = \mathrm{id}$. La réciproque est claire.

(b) Soit $N \trianglelefteq G$. La \omterm{ex:b3:representations:examples}{représentation régulière} de
$G/N$ est fidèle~; la décomposer en \omterm{def:b3:representations:rep}{irréductibles} de $G/N$ et
les relever à $G$ (\cref{prop:b3:representations:onedim}(b))~:
\omterm{def:b3:representations:character}{caractères} \omterm{def:b3:representations:rep}{irréductibles} $\chi_{i_1}, \dots$ de $G$ dont les
noyaux contiennent $N$ et dont le noyau \emph{commun} est
exactement l'image réciproque de $\{e\}$, i.e.\ $N$ (fidélité
sur le quotient). Donc $N = \bigcap_j \ker\chi_{i_j}$.

(c) Si $G$ est \omterm{def:b3:groups:simple}{simple}~: pour un \omterm{def:b3:representations:rep}{irréductible} non trivial $\chi$,
$\ker\chi \trianglelefteq G$ n'est pas $G$ (une \omterm{def:b3:representations:rep}{représentation
irréductible} triviale sur tout $G$ est le \omterm{def:b3:representations:character}{caractère} trivial),
donc $\ker\chi = \{e\}$. Réciproquement, supposons tous les
noyaux non triviaux triviaux, et soit $N \trianglelefteq G$ avec
$N \neq G$. Dans l'expression de $N$ en (b) comme intersection
de noyaux, un \omterm{def:b3:representations:character}{caractère} impliqué est non trivial (si tous
étaient triviaux, l'intersection serait $G$), et son noyau est
$\{e\}$~: $N = \{e\}$. Donc les seuls sous-groupes normaux sont
$\{e\}$ et $G$.
\end{solution}

\begin{solution}{exo:b3:representations:7}
Ordre $8$ non abélien~: le nombre de \omterm{def:b3:representations:character}{caractères} de \omterm{def:b3:galois:extension}{degré} $1$ est
$[G : D(G)]$, un diviseur propre de $8$ (non abélien~: $D(G)
\neq \{e\}$), et $\sum n_i^2 = 8$. Avec $k$ uns et les \omterm{def:b3:galois:extension}{degrés}
restants $\geq 2$~: $8 - k \equiv 0$ avec des carrés $\geq 4$,
et $k \mid 8$, $k < 8$. $k = 4$~: un \omterm{def:b3:galois:extension}{degré} $2$ --- cohérent. $k
= 2$~: $6$ restants, pas une somme de carrés $\geq 4$. $k =
1$~: impossible, car $k = [G : D(G)] \geq 2$ --- $G/D(G)$ est
un $2$-groupe abélien non trivial, car $G$ est un $2$-groupe
avec $D(G) \neq G$ par résolubilité des $p$-groupes
(\cref{ex:b3:groups:solvableexamples}). Donc le motif est
$(1,1,1,1,2)$. Fidélité de $\chi_5$~: les quatre \omterm{def:b3:representations:character}{caractères} de
\omterm{def:b3:galois:extension}{degré} $1$ contiennent tous $D(G)$ dans leurs noyaux~; si
$\ker\chi_5 \supseteq N \neq \{e\}$ pour un \omterm{def:b3:groups:normal}{normal} minimal $N$
($N \subseteq D(G)$ ou non --- prendre $N \subseteq \ker\chi_5$
non trivial), alors $N$ serait dans les cinq noyaux, dont
l'intersection est triviale (la \omterm{ex:b3:representations:examples}{représentation régulière} est
fidèle)~: contradiction. Ordre $p^3$ non abélien~: $Z(G)$ a
ordre $p$ (ordre $p^2$ rendrait $G/Z$ cyclique, $G$ abélien),
$G/Z(G)$ d'ordre $p^2$ est abélien, donc $D(G) \subseteq Z(G)$,
et $D(G) \ne \{e\}$~: $D(G) = Z(G)$, donnant $p^2$ \omterm{def:b3:representations:character}{caractères}
de \omterm{def:b3:galois:extension}{degré} $1$. Les \omterm{def:b3:galois:extension}{degrés} restants vérifient $\sum n_i^2 = p^3 -
p^2$ avec chaque $n_i > 1$ divisant $\abs G$
(\cref{pb:b3:representations:1}, question 8) d'où $n_i \in
\{p\}$ ($n_i = p^2$ dépasserait~: $p^4 > p^3 - p^2$)~: exactement
$p - 1$ \omterm{def:b3:representations:character}{caractères} de \omterm{def:b3:galois:extension}{degré} $p$.
\end{solution}

\begin{solution}{exo:b3:representations:8}
Définir, pour des \omterm{def:b3:representations:rep}{représentations} $\rho, \sigma$ de $G, H$ dans
$V, W$, la \omterm{def:b3:representations:rep}{représentation} $\rho \boxtimes \sigma$ de $G \times
H$ sur $V \otimes W$ --- concrètement, sur les matrices~:
$(\rho\boxtimes \sigma)(g,h)$ est le produit de Kronecker
$\rho(g)\otimes \sigma(h)$, dont la trace est
$\operatorname{tr}\rho(g)\operatorname{tr}\sigma(h) =
\chi(g)\psi(h)$ (le produit de Kronecker de matrices $A \otimes
B$ a pour trace $\operatorname{tr}A\operatorname{tr}B$~: sa
diagonale est $a_{kk}b_{ll}$). Donc $\chi\psi$ est un \omterm{def:b3:representations:character}{caractère},
et
\[
\langle\chi\psi, \chi'\psi'\rangle_{G\times H}
= \frac{1}{\abs G\abs H}\sum_{g,h}
\overline{\chi(g)\psi(h)}\,\chi'(g)\psi'(h)
= \langle\chi,\chi'\rangle_G\,\langle\psi,\psi'\rangle_H
= \delta_{\chi\chi'}\delta_{\psi\psi'}.
\]
En particulier $\langle\chi\psi,\chi\psi\rangle = 1$~: chaque
$\chi\psi$ est \omterm{def:b3:representations:rep}{irréductible}
(\cref{cor:b3:representations:multiplicity}). Ce sont $r_Gr_H$
\omterm{def:b3:representations:character}{caractères} \omterm{def:b3:representations:rep}{irréductibles} distincts~; les classes de $G\times H$
sont les produits de classes ($(g,h) \sim (g',h')$
composante par composante), donc il y en a $r_Gr_H$~: la liste
est complète (\cref{thm:b3:representations:numberirr}). Pour
$(\Z/2\Z)^2$~: les quatre \omterm{def:b3:representations:character}{caractères} de signe
$(\pm1)\otimes(\pm1)$ --- la table du bloc en haut à gauche de
\cref{exo:b3:representations:4}. Un groupe abélien fini est un
produit de groupes cycliques
(\cref{cor:b3:modules:abelian})~; ses \omterm{def:b3:representations:character}{caractères} \omterm{def:b3:representations:rep}{irréductibles}
sont des produits des cycliques
(\cref{exo:b3:representations:1})~: tous de \omterm{def:b3:galois:extension}{degré} $1$.
\end{solution}

\begin{solution}{exo:b3:representations:9}
(a) $D(A_5) = A_5$ (\omterm{def:b3:groups:simple}{simple} non abélien), donc le seul \omterm{def:b3:representations:character}{caractère}
de \omterm{def:b3:galois:extension}{degré} $1$ est $\mathbf 1$
(\cref{prop:b3:representations:onedim}). Il faut $n_2^2 + n_3^2
+ n_4^2 + n_5^2 = 59$ avec chaque $n_i \geq 2$~; en testant les
carrés $4, 9, 16, 25, 36, 49$~: le seul multiensemble qui
marche est $\{9, 9, 16, 25\}$~: avec plus grand carré $49$, le
reste $10$ n'est pas une somme de trois carrés $\geq 4$~; avec
$36$, le reste $23$ non plus ($16 + 4 + 4 = 24$, $9 + 9 + 4 =
22$)~; avec plus grand $25$, on vérifie $25 + 16 + 9 + 9 = 59$
marche et les variantes $25 + 25$, $25 + 16 + 16$, $25 + 16 +
4$ échouent~; avec plus grand $16$~: $16\cdot3 = 48 < 59 - 4$.
\omterm{def:b3:galois:extension}{Degrés}~: $1, 3, 3, 4, 5$.

(b) Permutation sur $5$ points~: points fixes $(5, 1, 2, 0,
0)$, donc $\chi_4 = (4, 0, 1, -1, -1)$ avec
$\langle\chi_4,\chi_4\rangle = \frac{16 + 0 + 20 + 12 + 12}{60}
= 1$~: \omterm{def:b3:representations:rep}{irréductible}. \omterm{def:b3:groups:action}{Action} sur les six $5$-Sylow~: une
involution en fixe exactement $2$ (les \omterm{ex:b3:groups:actions}{normalisateurs} sont
diédraux d'ordre $10$, chacun contenant $5$ involutions~: $30$
incidences pour $15$ involutions), un $3$-élément en fixe $0$
(pas d'ordre $3$ dans $D_5$), un $5$-élément en fixe
exactement $1$ (il appartient à un unique Sylow)~: \omterm{def:b3:representations:character}{caractère} de
permutation $(6, 2, 0, 1, 1)$, et $\chi_5 = (5, 1, -1, 0, 0)$
avec norme $\frac{25 + 15 + 20 + 0 + 0}{60} = 1$~:
\omterm{def:b3:representations:rep}{irréductible}. Deux lignes $(3, a, b, c, d)$, $(3, a', b', c',
d')$ restent. Normes de colonnes
(\cref{cor:b3:representations:column})~: sur la classe de
$(1\,2)(3\,4)$, $\abs{Z} = 4$~: $1 + 0 + 1 + a^2 + a'^2 = 4$~;
colonne contre $e$~: $1\cdot1 + 4\cdot 0 + 5\cdot1 + 3(a + a')
= 0$~: $a + a' = -2$, $a^2 + a'^2 = 2$~: $a = a' = -1$. Sur les
$3$-cycles, $\abs Z = 3$~: $1 + 1 + 1 + b^2 + b'^2 = 3$~: $b =
b' = 0$. Sur chaque $5$-classe, $\abs Z = 5$~: la colonne
contre $e$ lit $1\cdot 1 + 4\cdot(-1) + 5\cdot 0 + 3(c + c') =
0$, donc $c + c' = 1$~; et $c^2 + c'^2 = 3$~: $\{c, c'\} =
\bigl\{\frac{1+\sqrt5}2, \frac{1-\sqrt5}2\bigr\}$ --- le nombre
d'or et son conjugué~; la seconde $5$-classe porte les valeurs
échangées (les deux lignes doivent être orthogonales).

(c) Dans la table finie, aucune entrée d'une ligne non triviale
n'égale son \omterm{def:b3:galois:extension}{degré} hors de la première colonne~: chaque noyau
$\{g : \chi_i(g) = n_i\}$ est trivial. Par
\cref{exo:b3:representations:6}(c), $A_5$ est \omterm{def:b3:groups:simple}{simple}.
\end{solution}

\begin{solution}{exo:b3:representations:10}
$\rho(z)$ commute avec tout $\rho(g)$ ($z$ est central), i.e.\
$\rho(z) \in \operatorname{End}_G(V) = \C\,\mathrm{id}$
(Schur)~: $\rho(z) = \lambda_z\,\mathrm{id}$, et $z \mapsto
\lambda_z$ est multiplicatif~: un morphisme $Z(G) \to
\C^\times$. Alors $\chi(z) = n\lambda_z$ avec $\abs{\lambda_z}
= 1$ (racine de l'unité)~: $\abs{\chi(z)} = n$. Si $\rho$ est
fidèle, $\lambda$ est injectif sur $Z(G)$ ($\rho(z) =
\mathrm{id} \iff \lambda_z = 1$), donc $Z(G)$ s'injecte dans
$\C^\times$~; un sous-groupe fini du groupe multiplicatif d'un
corps est cyclique (\cref{thm:b3:galois:cyclic}).
\end{solution}

\begin{solution}{exo:b3:representations:11}
(a) Équivariance~: $\rho(h)p_i\rho(h)^{-1}$ réindexe la somme
($g \mapsto hgh^{-1}$, et $\chi_i$ est une \omterm{def:b3:representations:character}{fonction de
classe})~: $p_i$ commute avec l'\omterm{def:b3:groups:action}{action}. Sur une \omterm{def:b3:representations:rep}{irréductible} $W$
de \omterm{def:b3:representations:character}{caractère} $\chi_j$, Schur fait de la restriction un scalaire
$\lambda\,\mathrm{id}_W$~; en prenant les traces,
\[
\lambda\,n_j = \frac{n_i}{\abs G}\sum_g
\overline{\chi_i(g)}\,\chi_j(g) = n_i\,\langle\chi_i,
\chi_j\rangle = n_i\,\delta_{ij}
\]
(première orthogonalité)~: $\lambda = \delta_{ij}$.

(b) Décomposer $V$ en \omterm{def:b3:representations:rep}{irréductibles} (Maschke)~: $p_i$ agit
comme l'identité sur les facteurs de \omterm{def:b3:representations:character}{caractère} $\chi_i$ et
comme $0$ sur tous les autres, donc $p_i$ est la projection sur
leur somme $V_i$ le long de la somme du reste~; l'image $V_i$
ne dépend pas de la décomposition choisie (c'est l'ensemble
des vecteurs fixés par $p_i$, défini sans choix). $\sum_ip_i$
agit comme l'identité sur chaque facteur \omterm{def:b3:representations:rep}{irréductible}~: c'est
$\mathrm{id}_V$. Le scindage plus fin de $V_i \cong
W_i^{\oplus m_i}$ implique le choix d'une base de
$\operatorname{Hom}_G(W_i, V)$~: il n'est pas canonique.

(c) Pour $\varepsilon$ (\omterm{def:b3:galois:extension}{degré} $1$)~: $p_\varepsilon =
\frac1{6}\sum_{g}\varepsilon(g)\,\rho(g)$, i.e.\ dans l'algèbre
du groupe
\[
p_\varepsilon = \tfrac16\bigl(e - (1\,2) - (1\,3) - (2\,3)
+ (1\,2\,3) + (1\,3\,2)\bigr) .
\]
En élevant au carré~: le coefficient de $g$ dans
$p_\varepsilon^2$ est
$\frac1{36}\sum_{h}\varepsilon(h)\varepsilon(h^{-1}g) =
\frac1{36}\,\varepsilon(g)\sum_h\varepsilon(h)^2 =
\frac{6}{36}\varepsilon(g)$~: $p_\varepsilon^2 =
p_\varepsilon$. (Son image dans la \omterm{ex:b3:representations:examples}{représentation régulière} est
la droite engendrée par $\sum_g\varepsilon(g)e_g$~: la
\omterm{def:b3:representations:rep}{représentation} signature apparaît avec multiplicités $1$, comme
le demande la théorie générale.)
\end{solution}

\begin{solution}{exo:b3:representations:12}
(a) Ordonner les classes $e$ [1], transpositions [6],
$3$-cycles [8], $4$-cycles [6], doubles transpositions [3]. Le
second \omterm{def:b3:representations:character}{caractère} linéaire est $\chi_2 = \varepsilon$ avec
valeurs $1, -1, 1, -1, 1$. Pour le \omterm{def:b3:representations:character}{caractère} de \omterm{def:b3:galois:extension}{degré} $2$
$\chi_3$, l'orthogonalité de chaque colonne avec la colonne de
l'identité ($\sum_in_i\chi_i(g) = 0$ pour $g \neq e$) donne, sur
les transpositions~: $1 - 1 + 2\chi_3 + 3(\chi_4 + \chi_5) =
0$~; le tour de signature $\chi_5 = \varepsilon\chi_4$ (un
\omterm{def:b3:representations:character}{caractère} de \omterm{def:b3:galois:extension}{degré} $3$ multiplié par un linéaire reste
\omterm{def:b3:representations:rep}{irréductible} --- même norme) fait s'annuler $\chi_4 + \chi_5$
sur les classes impaires~: $\chi_3 = 0$ là. Sur les $3$-cycles~:
$1 + 1 + 2\chi_3(c_3) + 3(\chi_4 + \chi_5)(c_3) = 0$ avec
$\chi_5 = \chi_4$ sur les classes paires~; la colonne de $c_3$
avec elle-même donne des données en $\abs{\chi_3}^2$~; en
résolvant le petit système (utiliser aussi l'orthogonalité de
ligne de $\chi_3$ avec $\mathbf 1$ et $\varepsilon$)~: $\chi_3
= (2, 0, -1, 0, 2)$, puis $\chi_4 = (3, 1, 0, -1, -1)$ et
$\chi_5 = \varepsilon\chi_4 = (3, -1, 0, 1, -1)$. La table
complète~:
\begin{center}
\begin{tabular}{c|ccccc}
 & $e$ & $6\,t$ & $8\,c_3$ & $6\,c_4$ & $3\,v$\\
\hline
$\chi_1$ & $1$ & $1$ & $1$ & $1$ & $1$\\
$\chi_2$ & $1$ & $-1$ & $1$ & $-1$ & $1$\\
$\chi_3$ & $2$ & $0$ & $-1$ & $0$ & $2$\\
$\chi_4$ & $3$ & $1$ & $0$ & $-1$ & $-1$\\
$\chi_5$ & $3$ & $-1$ & $0$ & $1$ & $-1$\\
\end{tabular}
\end{center}
(Toutes les lignes ont norme $1$~; toutes les colonnes sont
orthogonales~: les vérifications passent.)

(b) Les données de classes $1, 6, 8, 6, 3$ avec ces \omterm{def:b3:galois:extension}{degrés} sont
celles de $S_4$ (types de cycles $e$, $2$, $3$, $4$, $2{+}2$).
Noyaux~: $\ker\chi_2 = \{g : \varepsilon(g) = 1\} = A_4$
(classes $e, c_3, v$~: $1 + 8 + 3 = 12$, indice $2$)~;
$\ker\chi_3 = \{g : \chi_3(g) = 2\}$ = classes $e, v$~: le
groupe de Klein $V_4$, d'ordre $4$, \omterm{def:b3:groups:normal}{normal}. $\chi_4, \chi_5$
sont fidèles ($\chi_i(g) = n_i$ seulement en $e$).
Intersections de noyaux~: $\{e\}$, $V_4$, $A_4$, $G$ --- par
\cref{exo:b3:representations:6}(b) ce sont \emph{tous} les
sous-groupes normaux de $S_4$.

(c) Les \omterm{def:b3:representations:character}{caractères} avec $V_4 \subseteq \ker$ sont $\chi_1,
\chi_2, \chi_3$~: ils se factorisent par $G/V_4$, un groupe
d'ordre $6$ possédant des \omterm{def:b3:galois:extension}{degrés} \omterm{def:b3:representations:rep}{irréductibles} $1, 1, 2$ ---
la table de $S_3$. Comme la table du quotient \emph{est} un
invariant complet parmi les deux groupes d'ordre $6$
($\Z/6\Z$ aurait six \omterm{def:b3:representations:character}{caractères} linéaires), $G/V_4 \cong
S_3$~: le quotient est visible dans la table comme le bloc des
lignes contenant $V_4$ dans leur noyau.
\end{solution}

\begin{solution}{pb:b3:representations:1}
\textbf{1.} Si $\alpha^n + c_{n-1}\alpha^{n-1} + \dots + c_0 =
0$ ($c_i \in \Z$), alors $\alpha^n \in \Z\text{-span}(1, \dots,
\alpha^{n-1})$, et par récurrence toute puissance l'est~:
$\Z[\alpha]$ est engendré par $1, \alpha, \dots, \alpha^{n-1}$.
Réciproquement soit $\Z[\alpha] = \Z g_1 + \dots + \Z g_m$.
Écrire $\alpha g_i = \sum_j m_{ij}g_j$ avec $M = (m_{ij}) \in
M_m(\Z)$~: le vecteur $g = (g_i)$ vérifie $(\alpha I - M)g =
0$~; en multipliant par la comatrice, $\det(\alpha I -
M)\,g_i = 0$ pour tout $i$, et comme $1 \in \Z[\alpha]$ est une
combinaison $\Z$-linéaire des $g_i$, $\det(\alpha I - M) =
0$~: $\alpha$ est racine du unitaire $\det(XI - M) \in \Z[X]$.

\textbf{2.} Si $\Z[\alpha]$ est engendré par les puissances de
$\alpha$ jusqu'à $n-1$ et $\Z[\beta]$ par celles de $\beta$
jusqu'à $m - 1$, alors $\Z[\alpha,\beta]$ est engendré par les
$nm$ produits $\alpha^i\beta^j$ (réduire tout monôme). Les
sous-anneaux $\Z[\alpha + \beta]$ et $\Z[\alpha\beta]$ sont des
$\Z$-sous-modules du $\Z$-module \omterm{def:b3:modules:free}{de type fini}
$\Z[\alpha, \beta]$, donc \omterm{def:b3:modules:free}{de type fini} ($\Z[\alpha,\beta]$,
engendré par $nm$ éléments, est image de $\Z^{nm}$~; un
\omterm{def:b3:modules:module}{sous-module} se tire en arrière en un \omterm{def:b3:modules:module}{sous-module} de $\Z^{nm}$,
\omterm{def:b3:modules:free}{libre} de rang $\leq nm$ par
\cref{thm:b3:modules:submodule}, et son image engendre). Par
la question 1, $\alpha + \beta$ et $\alpha\beta$ sont des
entiers \omterm{def:b3:galois:algebraic}{algébriques}.

\textbf{3.} Si $\frac pq$ (\omterm{def:b3:representations:rep}{irréductible}) est racine d'un
polynôme unitaire à coefficients entiers de \omterm{def:b3:galois:extension}{degré} $n$, le
théorème des racines rationnelles (dénominateurs~: $p^n =
-q(\cdots)$) donne $q \mid p^n$, donc $q = \pm1$.

\textbf{4.} $\chi(g)$ est une somme de racines de l'unité
(\cref{prop:b3:representations:charbasics}), chacune un entier
\omterm{def:b3:galois:algebraic}{algébrique} (racine de $X^N - 1$)~; conclure par la question 2.

\textbf{5.} Pour $h \in G$~: $\rho(h)S_C\rho(h)^{-1} =
\sum_{g\in C}\rho(hgh^{-1}) = S_C$ ($C$ est une classe). Par
Schur, $S_C = \omega_C\,\mathrm{id}$~; en prenant les traces,
$\abs C\,\chi(g_C) = \omega_C\, n$.

\textbf{6.} $S_CS_{C'} = \sum_{x \in C, y \in C'}\rho(xy) =
\sum_{z \in G} a(z)\rho(z)$ avec $a(z) = \#\{(x,y) \in C\times
C' : xy = z\}$. La conjugaison par $h$ bijectionne les solutions
pour $z$ avec celles pour $hzh^{-1}$~: $a$ est une \omterm{def:b3:representations:character}{fonction de
classe} à valeurs dans $\N$, donc $S_CS_{C'} =
\sum_{C''}a_{CC'C''}S_{C''}$. En substituant $S_C =
\omega_C\,\mathrm{id}$ partout et en identifiant les scalaires~:
$\omega_C\omega_{C'} = \sum_{C''}a_{CC'C''}\,\omega_{C''}$.

\textbf{7.} Soit $M$ le $\Z$-module engendré par $1$ et tous les
produits $\omega_{C_1}\cdots\omega_{C_k}$~; par la question 6
tout tel produit se réduit à une combinaison $\Z$-linéaire de
$1$ et des $\omega_C$~: $M$ est \omterm{def:b3:modules:free}{de type fini}, et $\omega_C M
\subseteq M$ pour chaque $C$. En particulier $\Z[\omega_C]
\subseteq M$ est \omterm{def:b3:modules:free}{de type fini} (\omterm{def:b3:modules:module}{sous-module}, comme en question
2), et la question 1 fait de $\omega_C$ un entier \omterm{def:b3:galois:algebraic}{algébrique}.

\textbf{8.} Pour un \omterm{def:b3:representations:rep}{irréductible} $\chi$ de \omterm{def:b3:galois:extension}{degré} $n$~:
\[
\frac{\abs G}{n} = \frac{\abs G}{n}\langle\chi,\chi\rangle
= \frac1n \sum_{g}\chi(g)\overline{\chi(g)}
= \sum_{C}\frac{\abs C\,\chi(g_C)}{n}\,\overline{\chi(g_C)}
= \sum_C \omega_C\,\overline{\chi(g_C)} .
\]
Chaque $\overline{\chi(g_C)} = \chi(g_C^{-1})$ est un entier
\omterm{def:b3:galois:algebraic}{algébrique} (question 4), donc le second membre en est un
(questions 2, 7)~; il est rationnel, donc un entier (question
3)~: $n \mid \abs G$.

\textbf{9.} Bézout~: $u\abs C + vn = 1$ avec $u, v \in \Z$.
Alors
\[
\frac{\chi(g_C)}{n} = u\,\frac{\abs C\,\chi(g_C)}{n} +
v\,\chi(g_C) = u\,\omega_C + v\,\chi(g_C),
\]
un entier \omterm{def:b3:galois:algebraic}{algébrique}.

\textbf{10.} Supposer $0 < \abs{\chi(g_C)} < n$ et poser
$\alpha = \chi(g_C)/n$. Toutes les valeurs $\chi(g)$ vivent
dans $\Q(\zeta_N)$, $N = \abs G$ (sommes de racines $N$-ièmes
de l'unité). Pour $\sigma \in
\operatorname{Gal}(\Q(\zeta_N)/\Q)$~: $\sigma$ envoie racines
de l'unité sur racines de l'unité ($\sigma(\zeta^k) =
\zeta^{ak}$, \cref{thm:b3:galois:cyclotomicirred}), donc
$\sigma(\chi(g_C))$ est encore une somme de $n$ racines de
l'unité~: $\abs{\sigma(\alpha)} \leq 1$~; de plus
$\sigma(\alpha)$ est un entier \omterm{def:b3:galois:algebraic}{algébrique} (même \omterm{def:b3:galois:algebraic}{polynôme
minimal} qu'$\alpha$). Le produit $P = \prod_\sigma
\sigma(\alpha)$ est fixé par tout le \omterm{def:b3:galois:galois}{groupe de Galois}, donc
rationnel (\cref{thm:b3:galois:fundamental}(1)), et c'est un
entier \omterm{def:b3:galois:algebraic}{algébrique} avec
\[
0 < \abs P \leq \abs\alpha < 1
\]
(aucun facteur ne s'annule~: $\sigma(\alpha) = 0$ forcerait
$\alpha = 0$). Cela contredit la question 3. Donc $\chi(g_C) =
0$ ou $\abs{\chi(g_C)} = n$, et dans ce dernier cas $\rho(g_C)$
est scalaire (\cref{prop:b3:representations:charbasics}).

\textbf{11.} Orthogonalité des colonnes ($C \neq \{e\}$)~:
$\sum_i \chi_i(e)\overline{\chi_i(g_C)} = 0$, i.e.\ $1 +
\sum_{i \geq 2} n_i\overline{\chi_i(g_C)} = 0$. Si tout
$\chi_i$ non trivial avec $p \nmid n_i$ s'annulait en $g_C$,
alors en regroupant le reste par leur facteur $p$~:
\[
-\frac1p = \sum_{i \geq 2,\ p \mid n_i}
\frac{n_i}{p}\,\overline{\chi_i(g_C)},
\]
un entier \omterm{def:b3:galois:algebraic}{algébrique} --- contredisant la question 3. Donc un
certain $\chi_i$ non trivial a $p \nmid n_i$ et $\chi_i(g_C)
\neq 0$~; comme $\abs C = p^k$, $\gcd(\abs C, n_i) = 1$, et la
question 10 fait de $\rho_i(g_C)$ un scalaire. Or $G$ \omterm{def:b3:groups:simple}{simple} non
abélien~: $\chi_i$ est fidèle
(\cref{exo:b3:representations:6}(c)), et $Z_i = \{g :
\rho_i(g) \text{ scalaire}\}$ est un \omterm{def:b3:groups:normal}{sous-groupe normal}
(l'image réciproque sous $\rho_i$ des scalaires, qui forment un
\omterm{def:b3:groups:normal}{sous-groupe normal} --- voire central --- de l'image) contenant
$g_C \neq e$~: $Z_i = G$. Alors $\rho_i(G)$ est abélien et
fidèle, rendant $G$ abélien~: contradiction. \emph{Aucun \omterm{def:b3:groups:simple}{groupe
simple} non abélien n'a de \omterm{ex:b3:groups:actions}{classe de conjugaison} de taille
puissance de premier $> 1$.}

\textbf{12.} Soit $G$ \omterm{def:b3:groups:simple}{simple} d'ordre $p^aq^b$. Si $b = 0$ ($G$
un $p$-groupe)~: $Z(G) \neq \{e\}$
(\cref{thm:b3:groups:pfixed}) est \omterm{def:b3:groups:normal}{normal}, donc $Z(G) = G$~:
abélien. Sinon prendre $Q$ un $q$-Sylow et $z \in Z(Q)
\setminus\{e\}$ (\cref{thm:b3:groups:pfixed} encore). Alors $Q
\subseteq Z_G(z)$, donc la classe de $z$ a taille $[G :
Z_G(z)]$ divisant $[G : Q] = p^a$~: une puissance de $p$. Si la
taille est $1$, $z \in Z(G)$~: le \omterm{ex:b3:groups:actions}{centre} est un \omterm{def:b3:groups:normal}{sous-groupe
normal} non trivial, donc $Z(G) = G$, abélien. Si la taille est
$p^k > 1$~: la question 11 l'interdit pour $G$ \omterm{def:b3:groups:simple}{simple} non
abélien. Dans les deux cas un \omterm{def:b3:groups:simple}{groupe simple} d'ordre $p^aq^b$
est abélien ($\cong \Z/p\Z$).

\textbf{13.} Récurrence sur $\abs G$ ($\abs G = 1$~: \omterm{def:b3:groups:derived}{résoluble}).
Si $G$ est \omterm{def:b3:groups:simple}{simple}, la question 12 le rend abélien, donc
\omterm{def:b3:groups:derived}{résoluble}. Sinon prendre $N$ un \omterm{def:b3:groups:normal}{sous-groupe normal} non trivial
propre~: $\abs N$ et $\abs{G/N}$ sont encore de la forme
$p^{a'}q^{b'}$ et plus petits, donc $N$ et $G/N$ sont
\omterm{def:b3:groups:derived}{résolubles} par récurrence, et $G$ est \omterm{def:b3:groups:derived}{résoluble}
(\cref{prop:b3:groups:derived}). --- Avec trois premiers,
l'étape clé échoue~: l'indice d'un Sylow n'est plus une
puissance de premier, donc la classe d'un élément central d'un
Sylow n'a pas forcément taille puissance de premier. Et un
échec est inévitable~: $A_5$, d'ordre $2^2\cdot3\cdot5$, est
\omterm{def:b3:groups:simple}{simple} et non \omterm{def:b3:groups:derived}{résoluble}.

\textbf{14.} Les classes et tailles sont
\cref{exo:b3:groups:11}(a). La classe de $S_5$ de $c$ a taille
$24$~; si elle restait une seule classe de $A_5$, le
dénombrement orbite--stabilisateur donnerait $\abs{Z_{A_5}(c)}
= 60/24$, non entier --- concrètement, $Z_{S_5}(c) = \langle
c\rangle$ a ordre $5$ et vit dans $A_5$, donc la classe de
$A_5$ de $c$ a taille $60/5 = 12$~: la classe de $S_5$ se
scinde en deux ($c$ et $c^2$ sont conjugués dans $S_5$ par une
permutation impaire seulement).

\textbf{15.} Un seul \omterm{def:b3:representations:character}{caractère} linéaire~: une \omterm{def:b3:representations:rep}{représentation} de
\omterm{def:b3:galois:extension}{degré} $1$ se factorise par $G/D(G)$, et $D(A_5) = A_5$ ($A_5$
est \omterm{def:b3:groups:simple}{simple} non abélien, et $D(A_5)$ est \omterm{def:b3:groups:normal}{normal}, non trivial ---
$A_5$ n'est pas abélien). Donc $n_1 = 1$ et $n_2^2 + n_3^2 +
n_4^2 + n_5^2 = 59$ avec chaque $n_i \geq 2$. Carrés
disponibles~: $4, 9, 16, 25, 36, 49$. Une somme de quatre
d'entre eux égale à $59$~: le plus grand doit être $25$ ($36 +
4 + 4 + 9 = 53$ n'ajuste pas~: $36 + 16 + 4 + 4 = 60$, $36 + 9
+ 9 + 4 = 58$, $36 + 16 + 9 + 4 = 65$ --- aucune combinaison
avec $36$ ou $49$ ne marche), et $59 - 25 = 34 = 16 + 9 + 9$
(la seule façon~: $16 + 16 + 4 = 36$, $25 + 9 + 4 = 38$, $25 +
4 + 4 = 33$)~: \omterm{def:b3:galois:extension}{degrés} $1, 3, 3, 4, 5$.

\textbf{16.} $\langle\pi, \mathbf 1\rangle = \frac1{60}(5 +
15\cdot1 + 20\cdot2 + 0 + 0) = 1$ (une \omterm{def:b3:groups:action}{orbite} --- le
dénombrement de Burnside), et $\langle\pi, \pi\rangle =
\frac1{60}(25 + 15 + 80 + 0 + 0) = 2$~: $\pi$ contient le
\omterm{def:b3:representations:character}{caractère} trivial une fois, et son autre constituant est un
seul \omterm{def:b3:representations:rep}{irréductible}. Donc $\chi_4 = \pi - \mathbf 1$ est
\omterm{def:b3:representations:rep}{irréductible}, de \omterm{def:b3:galois:extension}{degré} $4$, avec valeurs $4, 0, 1, -1, -1$.

\textbf{17.} Sur les paires, un élément fixe $\{i,j\}$ ssi il
fixe ou échange $i, j$~: les comptages sont $10$ ($e$), $2$
($t = (1\,2)(3\,4)$ fixe $\{1,2\}, \{3,4\}$), $1$ ($(1\,2\,3)$
fixe $\{4,5\}$), $0, 0$. Alors $\langle\pi_{10},
\pi_{10}\rangle = \frac1{60}(100 + 15\cdot4 + 20\cdot1) = 3$~:
trois constituants \omterm{def:b3:representations:rep}{irréductibles}, chacun une fois. Et
$\langle\pi_{10}, \mathbf 1\rangle = \frac1{60}(10 + 30 + 20)
= 1$, $\langle\pi_{10}, \chi_4\rangle = \frac1{60}(40 + 0 +
20\cdot1\cdot1 + 0 + 0) = 1$~: donc $\pi_{10} = \mathbf 1 +
\chi_4 + \chi$ avec $\chi$ \omterm{def:b3:representations:rep}{irréductible} de \omterm{def:b3:galois:extension}{degré} $10 - 1 - 4 =
5$ et valeurs $\chi_5 = \pi_{10} - \mathbf 1 - \chi_4 = (5, 1,
-1, 0, 0)$.

\textbf{18.} Écrire $a, a'$ pour les valeurs de $\chi_2,
\chi_3$ en $t$, et $b, b'$ en $s = (1\,2\,3)$~; les quatre sont
réelles ($t$ et $s$ sont conjugués à leurs inverses). Colonne
$t$ contre colonne $e$~: $1 + 3a + 3a' + 4\cdot0 + 5\cdot1 =
0$, donc $a + a' = -2$~; colonne $t$ avec elle-même~: $1 + a^2
+ a'^2 + 0 + 1 = \frac{60}{15} = 4$, donc $a^2 + a'^2 = 2$~;
d'où $(a + a')^2 = 4 = 2 + 2aa'$ donne $aa' = 1$ et $a = a' =
-1$. Colonne $s$ contre $e$~: $1 + 3(b + b') + 4 - 5 = 0$
donne $b + b' = 0$~; colonne $s$ avec elle-même~: $1 + b^2 +
b'^2 + 1 + 1 = \frac{60}{20} = 3$ donne $b = b' = 0$.

\textbf{19.} Colonne $c$ contre colonne $e$~: $1 + 3(x + y) +
4(-1) + 5\cdot0 = 0$, donc $x + y = 1$. Colonne $c$ contre
colonne $c^2$ (classes distinctes, orthogonales)~: $1 + xy +
yx + 1 + 0 = 0$, donc $xy = -1$. Ainsi $x, y$ résolvent $T^2 -
T - 1 = 0$~: $\{x, y\} = \{\varphi, \bar\varphi\}$ avec
$\varphi = \frac{1 + \sqrt5}2$. Cohérence~: $x^2 + y^2 =
(x+y)^2 - 2xy = 3 = \frac{60}{12} - 2$ --- matchant l'identité
de colonne sur elle-même $1 + x^2 + y^2 + 1 + 0 = 5$. La
table~:
\begin{center}
\begin{tabular}{c|ccccc}
 & $e$ & $15\,t$ & $20\,s$ & $12\,c$ & $12\,c^2$\\
\hline
$\chi_1$ & $1$ & $1$ & $1$ & $1$ & $1$\\
$\chi_2$ & $3$ & $-1$ & $0$ & $\varphi$ & $\bar\varphi$\\
$\chi_3$ & $3$ & $-1$ & $0$ & $\bar\varphi$ & $\varphi$\\
$\chi_4$ & $4$ & $0$ & $1$ & $-1$ & $-1$\\
$\chi_5$ & $5$ & $1$ & $-1$ & $0$ & $0$\\
\end{tabular}
\end{center}

\textbf{20.} $\norm{\chi_2}^2 = \frac1{60}\bigl(9 + 15\cdot1 +
0 + 12\varphi^2 + 12\bar\varphi^2\bigr) = \frac{9 + 15 +
36}{60} = 1$, en utilisant $\varphi^2 + \bar\varphi^2 =
(\varphi + \bar\varphi)^2 - 2\varphi\bar\varphi = 1 + 2 = 3$.
De même $\langle\chi_2, \chi_3\rangle = \frac1{60}(9 + 15 + 0 +
12(2\varphi\bar\varphi)) = \frac{9 + 15 - 24}{60} = 0$.
Divisibilité de \omterm{thm:b3:galois:finitefields}{Frobenius}~: $1, 3, 3, 4, 5$ divisent tous
$60$. Les valeurs irrationnelles $\varphi, \bar\varphi$ sont
la différence visible avec $S_5$~: dans $S_5$ tout élément est
conjugué à tous les générateurs de son groupe cyclique de même
type de cycle --- en particulier $c \sim c^2$ --- forçant des
valeurs de \omterm{def:b3:representations:character}{caractères} rationnelles (voire entières)~; dans
$A_5$ le scindage des $5$-cycles ouvre la porte à $\Q(\sqrt5)$.

\textbf{21.} Un \omterm{def:b3:groups:normal}{sous-groupe normal} $N$ est une réunion de
classes, contient $e$, et $\abs N \mid 60$. Les sommes
candidates~: $1{+}15 = 16$, $1{+}20 = 21$, $1{+}12 = 13$,
$1{+}24 = 25$, $1{+}15{+}20 = 36$, $1{+}15{+}12 = 28$,
$1{+}15{+}24 = 40$, $1{+}20{+}12 = 33$, $1{+}20{+}24 = 45$,
$1{+}12{+}12 = 25$, $1{+}15{+}20{+}12 = 48$, $1{+}15{+}20{+}24
= 60 - 12 = \dots$ listant toutes les sous-sommes propres
contenant $1$~: aucune de $13, 16, 21, 25, 28, 33, 36, 40, 45,
48, \dots$ ne divise $60$ sauf $1$ lui-même~: $N = \{e\}$ ou
$A_5$. Simplicité, lue sur cinq nombres. Et le critère de la
question 11 est visible aussi~: les tailles de classes $15 =
3\cdot5$, $20 = 4\cdot5$, $12 = 4\cdot3$ sont toutes composées
de deux premiers --- aucune classe puissance de premier,
exactement comme le critère de Burnside l'exige d'un \omterm{def:b3:groups:simple}{groupe
simple}.

\textbf{22.} Une rotation de $\R^3$ d'angle $\theta$ a pour
valeurs propres $1, \eu^{\iu\theta}, \eu^{-\iu\theta}$~: trace
$1 + 2\cos\theta$. Pour $\theta = \frac{2\pi}5$~:
$2\cos\frac{2\pi}5 = \frac{\sqrt5 - 1}2$, donc la trace est $1
+ \frac{\sqrt5-1}2 = \frac{1+\sqrt5}2 = \varphi$. L'élément
non trivial $\tau$ de $\operatorname{Gal}(\Q(\sqrt5)/\Q)$
appliqué entrée par entrée à une \omterm{def:b3:representations:table}{table de caractères} envoie
\omterm{def:b3:representations:character}{caractères} sur \omterm{def:b3:representations:character}{caractères} (il commute avec l'algèbre
définissante~: $\tau\circ\chi$ est le \omterm{def:b3:representations:character}{caractère} de la
\omterm{def:b3:representations:rep}{représentation} obtenue en transportant les matrices par $\tau$
sur les entrées, ou abstraitement~: les relations
d'orthogonalité sont $\Q$-rationnelles, donc $\tau$ permute
leurs solutions)~; $\tau$ fixe $\chi_1, \chi_4, \chi_5$
(valeurs rationnelles) et doit donc échanger $\chi_2$ et
$\chi_3$~: le second \omterm{def:b3:representations:character}{caractère} de \omterm{def:b3:galois:extension}{degré} $3$ porte les valeurs
conjuguées, sans calculer de matrice. Géométriquement, les
deux \omterm{def:b3:representations:rep}{représentations} sont l'\omterm{def:b3:groups:action}{action} icosaédrale et sa
composée avec un automorphisme extérieur de $A_5$
(conjugaison par une transposition), qui échange les deux
classes de $5$-cycles.

\textbf{23.} Pour $z \in Z(G)$, $\rho(z)$ commute avec tout
$\rho(g)$, donc par le lemme de Schur $\rho(z) = \lambda\,
\mathrm{id}$~; comme $z$ a ordre fini, $\lambda$ est une
racine de l'unité, et $\abs{\chi(z)} = \abs\lambda\,n = n$.
Alors
\[
\abs G = \abs G\,\langle\chi, \chi\rangle
= \sum_{g \in G}\abs{\chi(g)}^2
\geq \sum_{z \in Z(G)}\abs{\chi(z)}^2
= \abs{Z(G)}\,n^2,
\]
i.e.\ $n^2 \leq [G : Z(G)]$. Un groupe non abélien d'ordre $8$
($D_4$ ou $Q_8$) a cinq \omterm{ex:b3:groups:actions}{classes de conjugaison}, donc cinq
\omterm{def:b3:galois:extension}{degrés} \omterm{def:b3:representations:rep}{irréductibles} avec $\sum n_i^2 = 8$~; la seule façon
d'écrire $8$ comme somme de cinq carrés $\geq 1$ est $1 + 1 +
1 + 1 + 4$~: \omterm{def:b3:galois:extension}{degrés} $1, 1, 1, 1, 2$. Les deux groupes ont un
\omterm{ex:b3:groups:actions}{centre} d'ordre $2$, et le \omterm{def:b3:representations:character}{caractère} de \omterm{def:b3:galois:extension}{degré} $2$ atteint
l'égalité~: $2^2 = 4 = [G : Z(G)]$ --- la borne est optimale.
Pour $A_5$, $Z = \{e\}$ et la borne lit $n^2 \leq 60$~:
satisfaite par $1, 3, 3, 4, 5$ avec de la marge ($25 \leq 60$),
comme il se doit car l'égalité forcerait (par la même chaîne)
$\chi$ à s'annuler hors du \omterm{ex:b3:groups:actions}{centre}.

\textbf{24.} $\rho(g)^m = \mathrm{id}$ pour $m$ l'ordre de
$g$, donc $\rho(g)$ est annihilé par $X^m - 1$, scindé à
racines \omterm{def:b3:groups:simple}{simples} sur $\C$~: diagonalisable, avec valeurs propres
$\lambda_1, \dots, \lambda_n$ racines de l'unité, dans une
base propre $(e_i)$. Les produits $e_ie_j$ ($i \leq j$)
forment une base propre de $\operatorname{Sym}^2V$ avec
valeurs propres $\lambda_i\lambda_j$, et $e_i \wedge e_j$ ($i
< j$) une de $\Lambda^2V$~; comme
\[
\sum_{i<j}\lambda_i\lambda_j
= \frac{(\sum_i\lambda_i)^2 - \sum_i\lambda_i^2}2
= \frac{\chi(g)^2 - \chi(g^2)}2,
\]
et la somme symétrique ajoute $\sum_i\lambda_i^2$ au lieu de
la soustraire, les deux formules suivent. Pour $\chi_2 = (3,
-1, 0, \varphi, \bar\varphi)$ sur les classes $(e,
(2,2)\text{-}, 3\text{-cycles}, c, c^2)$~: élever au carré
envoie les doubles transpositions sur $e$, les $3$-cycles sur
des $3$-cycles, la classe de $c$ sur celle de $c^2$ et
réciproquement ($c^{-1} \sim c$ dans $A_5$, donc $c^4 \sim
c$). D'où $\chi_2(g^2)$ lit $(3, 3, 0, \bar\varphi,
\varphi)$, et, en utilisant $\varphi^2 = \varphi + 1$,
$\bar\varphi = 1 - \varphi$~:
\[
\Lambda^2\chi_2 = (3, -1, 0, \varphi, \bar\varphi) = \chi_2,
\qquad
\operatorname{Sym}^2\chi_2 = (6, 2, 0, 1, 1) .
\]
En décomposant ce dernier, avec les tailles de classes $1, 15,
20, 12, 12$~: $\langle\cdot, \mathbf 1\rangle = \frac1{60}(6 +
15\cdot2 + 0 + 12 + 12) = 1$~; $\langle\cdot, \chi_5\rangle =
\frac1{60}(30 + 30) = 1$~; $\langle\cdot, \chi_4\rangle =
\frac1{60}(24 - 12 - 12) = 0$~; $\langle\cdot, \chi_2\rangle =
\frac1{60}(18 - 30 + 12(\varphi + \bar\varphi)) = 0$, et de
même pour $\chi_3$. Donc $\operatorname{Sym}^2\chi_2 =
\mathbf 1 + \chi_5$ (dimensions $6 = 1 + 5$) et
$\chi_2\otimes\chi_2 = \mathbf 1 + \chi_2 + \chi_5$
(dimensions $9 = 1 + 3 + 5$). Géométrie~: l'isomorphisme
équivariant $\Lambda^2\R^3 \to \R^3$, $u \wedge v \mapsto u
\times v$, est exactement $\Lambda^2\chi_2 = \chi_2$ pour un
groupe de rotations~; le facteur $\mathbf 1$ du carré
symétrique est la forme quadratique invariante $x^2 + y^2 +
z^2$, et $\chi_5$ vit sur l'espace de dimension cinq des
tenseurs symétriques de trace nulle (quadratiques
harmoniques).

\textbf{25.} Colonne de $c$ contre colonne de $c^2$~:
\[
1\cdot1 + \varphi\bar\varphi + \bar\varphi\varphi +
(-1)(-1) + 0 = 1 - 1 - 1 + 1 + 0 = 0,
\]
comme l'exige l'orthogonalité pour des classes distinctes
($\varphi \bar\varphi = -1$). Colonne de $c$ contre
elle-même~: $1 + \varphi^2 + \bar\varphi^2 + 1 + 0 = 1 + 3 +
1 = 5 = 60/12 = \abs{Z_{A_5}(c)}$. \omterm{def:b3:representations:character}{Caractère} régulier
$\sum_in_i\chi_i$ sur les quatre colonnes non identité~:
\[
1 - 3 - 3 + 0 + 5 = 0, \qquad
1 + 0 + 0 + 4 - 5 = 0,
\]
sur doubles transpositions et $3$-cycles, et sur la classe de
$c$ (celle de $c^2$ est son conjugué galoisien)~:
\[
1 + 3\varphi + 3\bar\varphi - 4 + 0 = 1 + 3 - 4 = 0,
\]
en utilisant $\varphi + \bar\varphi = 1$. La table passe
chaque audit~: c'est la \omterm{def:b3:representations:table}{table de caractères} de $A_5$.
\end{solution}
